% Сборка: pdflatex lab02.tex
\documentclass[a4paper,12pt]{article}

\usepackage{cmap}
\usepackage[T1,T2A]{fontenc}
\usepackage[utf8]{inputenc}
\usepackage[english,russian]{babel}

\usepackage{amsmath}
\usepackage[a4paper,margin=2.5cm]{geometry}
\usepackage{enumitem}
\usepackage{listings}

% Чёрно-белое оформление кода: комментарии курсивом, без цвета.
\lstdefinestyle{bw}{
  language=Python,
  basicstyle=\ttfamily\small,
  commentstyle=\itshape,
  numbers=none,
  frame=single,
  columns=fullflexible,
  keepspaces=true,
  showstringspaces=false,
  breaklines=true,
  upquote=true,
}
\lstset{style=bw}

\newcommand{\pts}[1]{\textbf{(#1)}}

\begin{document}

\begin{center}
  {\Large\bfseries Лабораторная работа 2\par}
  {\large Метод максимального правдоподобия\par}
\end{center}

\textbf{Как сдавать работу.}
\begin{enumerate}[itemsep=0.3em]
  \item Сохраните копию ноутбука на свой Диск («Файл $\to$ Сохранить копию на Диске»).
  \item Решите задания в своей копии.
  \item Перед отправкой выберите «Среда выполнения $\to$ Перезапустить и выполнить всё» и убедитесь, что выводы всех ячеек (числа, графики) видны.
  \item В конце ноутбука напишите 2--3 предложения вывода по практической части.
  \item Откройте доступ: «Поделиться $\to$ Все, у кого есть ссылка $\to$ Читатель».
  \item Пришлите преподавателю ссылку на свою копию.
\end{enumerate}

Всего за работу можно получить 10 баллов.

Обозначения: $X_i$~--- элементы случайной выборки (случайные величины), $x_i$~--- их реализация (наблюдаемые значения).

\section*{Теоретические вопросы}

\begin{enumerate}
  \item \pts{1 балл} Дайте определение функции правдоподобия для реализации $\mathbf{x}=\left\{x_i\right\}_{i=1}^n$ случайной выборки $\mathbf{X}=\left\{X_i\right\}_{i=1}^n$ из распределения с плотностью (или функцией вероятности) $p(x\mid\theta)$, где $\theta$~--- неизвестный параметр. Дайте определение оценки максимального правдоподобия (ОМП) $\hat\theta_n$.

  \item \pts{1 балл} Почему для нахождения ОМП удобнее максимизировать логарифм функции правдоподобия (log-likelihood), а не саму функцию правдоподобия?

  \item \pts{2 балла} Напишите функцию правдоподобия и логарифм функции правдоподобия для реализации $\mathbf{x}=\left\{x_i\right\}_{i=1}^n$ случайной выборки $\mathbf{X}=\left\{X_i\right\}_{i=1}^n$ из распределения Пуассона $\mathrm{Pois}(\lambda)$ с неизвестным параметром $\lambda>0$. Приравняв производную логарифма правдоподобия по $\lambda$ к нулю, найдите оценку максимального правдоподобия $\hat\lambda_n$. Убедитесь (по знаку второй производной), что это точка максимума.

  \item \pts{1 балл} Докажите, что оценка $\hat\lambda_n$ из вопроса 3 является несмещенной.

  \item \pts{1 балл} Докажите, что оценка $\hat\lambda_n$ из вопроса 3 является состоятельной.

\end{enumerate}

\section*{Практические задачи}

\begin{enumerate}
  \item \pts{1 балл} Выберите какое-нибудь значение параметра $\lambda$ распределения Пуассона и с помощью генератора \lstinline|np.random.default_rng(seed)| и его метода \lstinline|rng.poisson| сгенерируйте выборку объема $n=1000$. Постройте на одном графике эмпирические частоты значений выборки (столбчатая диаграмма) и теоретическую функцию вероятности $\mathrm{Pois}(\lambda)$ (\lstinline|scipy.stats.poisson.pmf|) и убедитесь, что эмпирические частоты приближают теоретические.

  \item \pts{1 балл} Реализуйте функцию, вычисляющую оценку максимального правдоподобия $\hat\lambda_n$ по реализации $\mathbf{x}=\left\{x_i\right\}_{i=1}^n$ случайной выборки $\mathbf{X}=\left\{X_i\right\}_{i=1}^n$ объема $n$ из распределения Пуассона:
\begin{lstlisting}
def hat_lambda_mle(x: np.ndarray) -> float:
    # your code
    return lam
\end{lstlisting}

  \item \pts{1 балл} Постройте график логарифма функции правдоподобия $\ell(\lambda)$ (из вопроса 3) от $\lambda$ для сгенерированной в задаче 1 выборки. Отметьте на графике вертикальной линией значение $\hat\lambda_n$, вычисленное функцией \lstinline|hat_lambda_mle|, и убедитесь, что оно соответствует максимуму графика.

  \item \pts{1 балл} Для выбранного в задаче 1 значения $\lambda$ постройте график зависимости оценки $\hat\lambda_n$ от объема выборки $n$ (по первым $n$ элементам выборки). Сделайте это для нескольких (например, 5) независимых выборок объема $1000$, изобразив их на одном графике, и отметьте на нем истинное значение $\lambda$. Убедитесь, что при увеличении $n$ оценки приближаются к истинному значению параметра, а разброс между выборками уменьшается.

\end{enumerate}

\section*{Вывод}

Напишите 2--3 предложения вывода по практической части: что вы наблюдали для эмпирических частот и функции вероятности, для графика логарифма правдоподобия и для сходимости оценки $\hat\lambda_n$ при росте $n$.

\end{document}
